% ---------------------------------------------------------------------------
% AAAI / IJCAI Conference Paper Template
%
% Usage:
%   1. Download official aaai24.sty from: https://aaai.org/authorkit24-2/
%   2. Place aaai24.sty and aaai24.bst in this directory
%   3. Compile: pdflatex template.tex && bibtex template && pdflatex template.tex x2
%
% Conference: AAAI (Association for the Advancement of Artificial Intelligence) / IJCAI
% Page limit: 7 pages (AAAI) + references (+ 2 extra pages for appendices in some years)
% Two-column format, camera-ready style
% ---------------------------------------------------------------------------

\documentclass[letterpaper]{article}

% --- UNCOMMENT when aaai24.sty is available ---
% \usepackage{aaai24}
% \usepackage{times}
% \usepackage{helvet}
% \usepackage{courier}
% \usepackage[hyphens]{url}
% \usepackage{graphicx}
% \urlstyle{rm}
% \def\UrlFont{\rm}

\usepackage[utf8]{inputenc}
\usepackage[T1]{fontenc}
\usepackage{graphicx}
\usepackage{amsmath,amssymb}
\usepackage{booktabs}
\usepackage{multirow}
\usepackage{times}
\usepackage[pagebackref,breaklinks,colorlinks,citecolor=blue]{hyperref}
\usepackage[margin=1in]{geometry}  % Remove when using aaai24.sty

% --- Paper ID ---
% \setAAAIID{12345}

\title{Your Paper Title Here}

\author{
  Author One${}^1$,
  Author Two${}^1$,
  Author Three${}^2$ \\
  \\
  ${}^1$Institution One, Country \\
  ${}^2$Institution Two, Country \\
  \texttt{\{email1, email2\}@institution.edu}
}

\begin{document}

\maketitle

% ===========================================================================
% ABSTRACT
% ===========================================================================
\begin{abstract}
  Your abstract here. AAAI papers span all AI subfields.
  Follow 2-6-2 structure: 2 background/gap, 6 method, 2 results/implication.
  Target: $\sim$200 words.

  AAAI emphasizes: broad AI relevance, practical impact, clear methodology.
  Ensure the abstract is accessible to a broad AI audience (not just your subfield).
\end{abstract}

% ===========================================================================
% 1. INTRODUCTION
% ===========================================================================
\section{Introduction}

% 5-paragraph structure, broad-AI-audience flavor
% AAAI papers benefit from: accessible motivation, clear contributions,
% and explicit connection to broader AI themes

Our main contributions are:
\begin{itemize}
  \item Contribution 1 -- [what + evidence].
  \item Contribution 2 -- [what + evidence].
  \item Contribution 3 -- [what + evidence].
\end{itemize}

% ===========================================================================
% 2. RELATED WORK
% ===========================================================================
\section{Related Work}

% Organize by themes with clear differentiation from your method
% AAAI reviewers appreciate explicit positioning within the AI landscape

% ===========================================================================
% 3. METHOD
% ===========================================================================
\section{Method}
\label{sec:method}

% Opening paragraph (3-5 sentences): overall framework → what each innovation
% solves → how they collaborate. No "Overview" subsection.
% Reference the architecture diagram.

\begin{figure}[t]
  \centering
  \fbox{\parbox{0.8\linewidth}{\centering
  \vspace{2cm}
  {\large [Architecture / Algorithm Diagram]}
  \vspace{2cm}}}
  \caption{Overview of the proposed method.}
  \label{fig:method}
\end{figure}

\subsection{<Innovation 1 Name>}
\label{sec:method-innov1}

% Motivation → Design → Formula/Mechanism → Effect

\subsection{<Innovation 2 Name>}
\label{sec:method-innov2}

% Motivation → Design → Formula/Mechanism → Effect

\subsection{<Innovation 3 Name>}
\label{sec:method-innov3}

% Motivation → Design → Formula/Mechanism → Effect

\subsection{Implementation Details}
\label{sec:method-impl}

% Optimizer, LR schedule, batch size, data augmentation, etc.

% ===========================================================================
% 4. EXPERIMENTS
% ===========================================================================
\section{Experiments}
\label{sec:experiments}

\subsection{Datasets and Metrics}
\label{sec:exp-data}

\subsection{Experimental Setup}
\label{sec:exp-setup}

\subsection{Comparison with State-of-the-Art and Classical Methods}
\label{sec:exp-main}

\subsection{Ablation Study}
\label{sec:exp-ablation}

% Numbered list 1), 2), 3)... analyzing each component / design choice

\subsection{Visualization}
\label{sec:exp-visual}

% ===========================================================================
% 5. CONCLUSION
% ===========================================================================
\section{Conclusion}
\label{sec:conclusion}

% Paragraph 1: Full summary (problem → method → key results with numbers)
% Paragraph 2: Limitations + future work (specific limitations + actionable directions)

% ===========================================================================
% ACKNOWLEDGMENTS
% ===========================================================================
\section*{Acknowledgments}

% ===========================================================================
% REFERENCES
% ===========================================================================
{
  \small
  \bibliographystyle{aaai24}
  \bibliography{references}
}

% ===========================================================================
% APPENDIX (optional -- check current AAAI CFP for appendix policy)
% ===========================================================================
% \appendix
% \section{Appendix}

\end{document}
